\subsection*{参数空间里的高度是什么}
网络参数很多，每一组参数都给一组预测。用损失函数评价预测误差，就可以把参数想象成位置、损失想象成高度。学习过程希望移动到较低的位置；虽然高维地形画不出来，局部变化仍能用导数描述。

对小位移 $d$，损失的一阶变化约为 $\nabla L\cdot d$。若固定步长大小，这个点积在反梯度方向最小。因此选择 $d=-\eta\nabla L$，预测变化为 $-\eta\|\nabla L\|^2$。这个理由依赖局部近似，不是对任何大步都有效的保证。

\subsection*{用一个可精确求解的地形检验直觉}
在 $L(w)=w^2$ 中，梯度为 $2w$，更新后 $w_+=(1-2\eta)w$。损失每步乘以 $(1-2\eta)^2$。现在能精确看见学习率的作用：小步缓慢靠近，较大但合适的步长可能左右摆动着靠近，过大则越走越远。

在交互实验中先选 $\eta=0.2$，再选 $0.8$，最后选 $1.1$。运行前分别预测序列的符号与绝对值变化；动画给出的应是公式的结果，而不是“训练总会下降”的预设故事。

真实网络的地形更复杂，可能有平坦方向、鞍点和噪声梯度。即使训练损失下降，也不直接保证对新数据有效。因此本节的二次模型用于理解更新机制与步长条件，不能被扩展成所有模型全局收敛的证明。
